% Procedencia: integral 2249, cap. 25, pp. 619--640, especialmente §25.13.
% Requiere la incorporación material de exc:k-direccion, exc:leech y el núcleo común.
\section{Dimensional screens and coordinated state reduction}
\label{iv:pantallas}
\begingroup
\setlength{\parskip}{2pt plus 1pt minus .5pt}
\setlength{\abovedisplayskip}{6pt plus 1pt minus 1pt}
\setlength{\belowdisplayskip}{6pt plus 1pt minus 1pt}
\setlength{\abovedisplayshortskip}{2pt plus 1pt}
\setlength{\belowdisplayshortskip}{4pt plus 1pt minus 1pt}

The ranks appearing in the exceptional prolongation and in
compactification belong to different objects. Relating them
requires retaining their presentation: generators, relations, inner
product, and selected direction. Dimension is then a
consequence of the channel module and its relations. Its
transport is formulated through maps acting on that module,
as well as on the state determining it.

\subsection{Channels, relations, and effective dimension}

Let \(\mathcal S_n\) be the set of surviving states of
depth \(n\), with compatible restrictions
\(r_{m,n}:\mathcal S_m\to\mathcal S_n\), \(m\geq n\).
The boundary datum \(\beta_n(x)\) preserves the information
required to decide compatibility of a prolongation:
sheet, orientation, quotient, memory, and incidence class according to
the readout considered. The restrictions satisfy
\[
 r_{\ell,n}=r_{m,n}r_{\ell,m},\qquad
 \beta_n r_{m,n}=b_{m,n}\beta_m.
\]
These equalities are the naturality conditions of the
presentation and retain the same data under restriction
in stages or directly.

\begin{definition}[Channel screen]
\label{iv:def:pantalla}
Fix a field \(\mathbb K_n\), and let \(M_n(x)\) be the space
generated by the prolongation channels and \(\mathcal R_n(x)\)
the subspace of relations. The screen is
\[
 \Sigma_n(x)=\bigl(\mathbb K_n,M_n(x),\beta_n(x),
                    \mathcal R_n(x)\bigr),
\]
and its effective dimension is
\begin{equation}
 d_{\Sigma_n}(x)=\dim_{\mathbb K_n}
                    \bigl(M_n(x)/\mathcal R_n(x)\bigr).
 \label{iv:eq:dimension-pantalla}
\end{equation}
If \(M_n(x)\cong\mathbb K_n^{g_n}\) and the columns of
\(R_n(x)\) generate the relations, then
\(d_{\Sigma_n}(x)=g_n-\operatorname{rank}R_n(x)\).
\end{definition}

\begin{proposition}[Invariance of presentation and transport]
\label{iv:prop:pantalla-transporte}
Dimension \eqref{iv:eq:dimension-pantalla} is preserved
under invertible changes of basis. For two screens over the same field
\(\mathbb K_n=\mathbb K_m=\mathbb K\), if the linear map
\(F_\gamma:M_n(x)\to M_m(T_\gamma x)\) sends
\(\mathcal R_n(x)\) into \(\mathcal R_m(T_\gamma x)\), it induces
a unique linear map between the screen quotients.
\end{proposition}
\begin{proof}
A change of generators and relations transforms \(R\) into
\(URV\), with \(U,V\) invertible, preserving its rank. For
descent, define
\(\overline F_\gamma([v])=[F_\gamma(v)]\). If two
representatives differ by a relation, linearity implies that their images differ
by a relation of the terminal screen; the map is
well defined. The quotient projection is surjective and
guarantees uniqueness.
\end{proof}

The dimensional change during transport is determined
by
\[
 \Delta d(\gamma;x)=d_{\Sigma_m}(T_\gamma x)
                         -d_{\Sigma_n}(x).
\]
This expression retains the state and its channel presentation.
An isolated rank permits reconstruction of neither the relations nor
the map that produced it.

\subsection{Ternary puncturing and syndrome screen}

The code \(C_W=[12,6,6]_3\), constructed from the
exceptional incidence, produces another screen by
puncturing one coordinate. Denote that map by
\(p:C_W\to\FF_3^{11}\) and its image by
\(G_{11}=p(C_W)\).

\begin{proposition}[Punctured code and perfect partition]
\label{iv:prop:codigo-once}
The map \(p\) is injective. The code \(G_{11}\) has
parameters \([11,6,5]_3\), and its Hamming balls of radius
\(2\) partition \(\FF_3^{11}\).
\end{proposition}
\begin{proof}
A vector in the kernel of \(p\) would have support contained in
the sole removed coordinate. The minimum distance \(6\)
of \(C_W\) excludes every nonzero vector of that form.
Thus the dimension remains \(6\). Puncturing decreases
weight by at most one, so the minimum distance
of the image is at least \(5\). The balls of radius \(2\)
are therefore disjoint. Each contains
\[
 1+2\binom{11}{1}+2^2\binom{11}{2}
 =1+22+220=243
\]
vectors. Since \(|G_{11}|=3^6\) and
\(3^6\cdot243=3^{11}\), the balls cover the ambient space.
To verify that the distance is exactly \(5\),
take a vector of weight \(3\). Its decoding ball
cannot have centre zero and has a centre of weight at
most \(3+2=5\); that centre must have weight at least
\(5\). A word of weight \(5\) is thereby obtained.
\end{proof}

The integer \(243=3^{11-6}\) is here the number of classes
of the syndrome quotient \(\FF_3^{11}/G_{11}\). Its
function is fixed by the puncturing and distance of the
code; equality in cardinality with another fibre does not replace
a map between them. This combinatorial screen and the
following real screen retain their own operations.

\subsection{The dodecaphase direction and rank reduction}

The construction in \ref{exc:k-direccion} starts from the record
\(K\), already obtained by APP--TRIT--TPK, and the
representation chart specified on \(\RR^{12}\). In that chart,
\begin{equation}
 V_{11}=\mathbf1^\perp,\qquad
 P_{11}=I_{12}-\frac1{12}J_{12},\qquad
 u_K=P_3P_{11}K.
 \label{iv:eq:p11-uk}
\end{equation}
Here \(J_{12}\) is the all-ones matrix and \(P_3\) is the
isotypic projector specified by the action of \(A_5\).
The finite sum defining \(P_3\), its representatives, and the
generation of \(K\) are retained in the cited section of
this same article. In particular,
\[
 \|u_K\|^2=\frac{6638585+2275584\sqrt5}{20}>0.
\]
The positive norm defines \(\ell_K=\RR u_K\), and the
projector
\begin{equation}
 P_{10}=P_{11}-\frac{u_Ku_K^{\mathsf T}}{\|u_K\|^2}
 \label{iv:eq:p10}
\end{equation}
has image \(V_{10}=V_{11}\cap u_K^\perp\).

\begin{proposition}[Two orthogonal reductions]
\label{iv:prop:12-11-10}
One has
\[
 \RR^{12}=\RR\mathbf1\mathbin{\oplus^\perp}V_{11},\qquad
 V_{11}=\ell_K\mathbin{\oplus^\perp}V_{10},\qquad
 \dim V_{10}=10.
\]
Moreover, \(P_{10}P_{11}=P_{10}\), \(P_{10}u_K=0\), and
\(P_{10}\) is an orthogonal projector.
\end{proposition}
\begin{proof}
Proposition \ref{prop:k-reduccion-diez}, proved before
the screens are constructed, establishes idempotence, symmetry, rank, and
the two operator identities for this same \(u_K\) and this same
\(P_{10}\). The first decomposition follows from
\(P_{11}=I-J_{12}/12\); the second is the orthogonal decomposition
of that proposition. The single proof of the
reduction is thus retained, and its two stages are made explicit here.
\end{proof}

The first reduction removes only the uniform mode.
The second removes the direction determined by \(K\) in
the fixed representation. If the chart is simultaneously changed
by an orthogonal matrix \(S\), the vector transforms as
\(u_K\mapsto Su_K\) and the projector as
\(P_{10}\mapsto SP_{10}S^{-1}\). This covariance preserves
the geometric operation although its coordinates change.

\subsection{Lorentzian lattice screen}

Let \(\Lambda=N_\alpha\) be the Leech lattice constructed
in \ref{exc:leech}. Consider the integral hyperbolic plane
\(U=\ZZ e_+\oplus\ZZ e_-\), with
\[
 \langle e_+,e_+\rangle=\langle e_-,e_-\rangle=0,
 \qquad\langle e_+,e_-\rangle=1.
\]
The lattice \(L_{26}=\Lambda\oplus U\) is even and
unimodular, of signature \((25,1)\). Each vector is written
\(y=\lambda+ae_++be_-\). Since
\(\langle y,e_+\rangle=b\),
\begin{equation}
 e_+^\perp=\Lambda\oplus\ZZ e_+,
 \qquad e_+^\perp/\ZZ e_+\cong\Lambda.
 \label{iv:eq:26-24}
\end{equation}
The quotient map is
\(q_{24}(\lambda+ae_+)=\lambda\); it is surjective and has
kernel \(\ZZ e_+\). The rank decreases by two through
an orthogonality restriction followed by an isotropic
quotient. This is therefore an operation of a different type from
\eqref{iv:eq:p10}.

\subsection{Joint screen morphism}

With the domain \(\mathscr D_0\) of \ref{iv:def:dualidad},
let
\[
 \mathscr G_K=\{(v_{11},P_{10}v_{11}):v_{11}\in V_{11}\},
\]
\[
 \mathscr S_{\mathrm{HMT}}=\RR^{12}\times\mathscr D_0\times e_+^\perp,
 \qquad
 \mathscr S_M^{\mathrm{alg}}=\mathscr G_K\times\mathscr D_0\times\Lambda.
\]
\begin{definition}[Algebraic map of coordinated screens]
\label{iv:def:rm-alg}
Define
\begin{equation}
 \mathcal R_M^{\mathrm{alg}}(v,d,y)
 =\bigl((P_{11}v,P_{10}v),d,q_{24}(y)\bigr).
 \label{iv:eq:rm-alg}
\end{equation}
\end{definition}

\begin{theorem}[Quotient isomorphism and intertwining of duality]
\label{iv:thm:pantallas-isomorfas}
\label{thm:k-pantallas-coordinadas}
The map \eqref{iv:eq:rm-alg} induces an isomorphism
\[
 \overline{\mathcal R}_M^{\mathrm{alg}}:
 (\RR^{12}/\RR\mathbf1)\times\mathscr D_0
 \times(e_+^\perp/\ZZ e_+)
 \longrightarrow\mathscr S_M^{\mathrm{alg}}.
\]
It intertwines the transformations acting by \(\mathsf T_0\)
on the factor \(\mathscr D_0\) and by the identity on the
other two factors. The form \(E_{R_0}(m,w)\) is preserved
under that duality.
\end{theorem}
\begin{proof}
The map \([v]\mapsto P_{11}v\) is an isomorphism
from \(\RR^{12}/\RR\mathbf1\) onto \(V_{11}\).
Since \(P_{10}P_{11}=P_{10}\), the map
\([v]\mapsto(P_{11}v,P_{10}v)\) identifies that quotient
with \(\mathscr G_K\); its inverse takes the first component
and its class. By \eqref{iv:eq:26-24}, \(q_{24}\) induces the
lattice isomorphism of the third factor. The second is
transported by the identity. The product of the three
maps proves the isomorphism.

The projectors and \(q_{24}\) do not modify \(d\), while
duality acts only on \(d\). The two
compositions are therefore identical. Preservation of the form
is Theorem \ref{iv:thm:dualidad-escalar}.
\end{proof}

The graph \(\mathscr G_K\) retains the preceding coordinate
\(v_{11}\) together with its projection. The theorem proves
equivalence of the quotient presentations; the isolated reduction
\(V_{11}\to V_{10}\) retains its kernel
\(\ell_K\). This precision establishes what information
the isomorphism retains and what information is removed by each
projection composing it.

\begin{corollary}[Screens over the elliptic collection]
The same morphism restricts to an isomorphism
\[
 (\RR^{12}/\RR\mathbf1)\times\mathscr D_{\mathrm{el}}
 \times(e_+^\perp/\ZZ e_+)
 \ \cong\ \mathscr G_K\times\mathscr D_{\mathrm{el}}\times\Lambda,
\]
and preserves the intertwining of duality.
\end{corollary}
\begin{proof}
The morphism and its inverse preserve the second factor exactly.
Corollary \ref{iv:cor:radio-especializado} proves that duality
preserves \(\mathscr D_{\mathrm{el}}\). Thus both the bijection
and the already proved intertwining square restrict to it.
\end{proof}

\subsection{Coordination with the exceptional realization}

On the domain \(\mathcal X_{\mathrm{HMT}}\), with phase, orientation,
quotient, boundary, incidence, grading, and memory, the readouts
\(\varepsilon_{\mathrm{exc}}\) and \(\varepsilon_{\mathrm{scr}}\) extract
the exceptional and screen data. The corridor constructs the graded
algebra \(\mathfrak V(x)\cong V^\natural\), with its vertex product
and automorphisms; \(\mathfrak S_M(x)\) assembles the preceding screens, quotients,
and operators. The joint assignment
\(x\mapsto(\mathfrak V(x),\mathfrak S_M(x))\) therefore has an
algebra as its first component and a collection of spaces and maps as its second.
The Monster action belongs to the first; the reduction directed by \(K\)
and compact duality belong to the second. The two readouts establish their common
provenance, and each construction retains its type.
Theorem~\ref{iv:thm:pantallas-isomorfas} gives the isomorphism
factor by factor; its proof identifies the inverses, kernels
and operations preserved by each component.
\par\endgroup
